\documentclass[12pt]{article}
\usepackage[utf8]{inputenc}
\usepackage{amsmath}
\usepackage{amssymb}
\usepackage{algorithm}
\usepackage{algpseudocode}
\usepackage{geometry}
\usepackage{booktabs}
\usepackage{tabularx}
\usepackage{hyperref}
\usepackage{xcolor}
\usepackage{titlesec}
\geometry{a4paper, margin=1in}

\definecolor{ppacblue}{RGB}{30, 80, 160}
\definecolor{roadgreen}{RGB}{20, 130, 60}

\hypersetup{colorlinks=true, linkcolor=ppacblue, urlcolor=ppacblue}

\title{\textbf{Police Patrol Area Covering Models}\\[0.4em]
       \large A Comparative Study: The Classical PPAC Model\\
       vs.\ a Road-Network Distance Extension}
\author{Graduation Project --- Determining Police Patrol Areas\\
        Based on Curtin, Hayslett-McCall \& Qiu (2010)}
\date{\today}

\begin{document}

\maketitle
\tableofcontents
\newpage

% =============================================================================
\section{Introduction}
% =============================================================================

Virtually every police department divides its jurisdiction into hierarchical
patrol geographies --- precincts, sectors, and beats.  Despite their operational
importance, these boundaries have historically been drawn by hand, with no
formal quantitative method for evaluating whether the resulting arrangement
is efficient \cite{curtin2010}.

Curtin, Hayslett-McCall, and Qiu (2010) introduced the \textbf{Police Patrol
Area Covering (PPAC) model}, which applies the Maximal Covering Location
Problem (MCLP) \cite{church1974} to optimally position patrol command centres
so as to maximise the number of weighted crime incidents reachable within an
acceptable service distance~$S$.  Applied to Dallas, TX, PPAC achieved an
18.9\,\% reduction in total response distance and an 18-percentage-point
improvement in incident coverage compared with the hand-drawn geography.

This document describes the PPAC model as formulated in the paper and then
presents an \textbf{extended, road-network-aware} version that replaces the
straight-line distance assumption of the original with shortest-path distances
computed over the OpenStreetMap (OSM) drivable road network.  A structured
side-by-side comparison of assumptions, formulations, and solution procedures
is given throughout.

% =============================================================================
\section{Shared Notation}
% =============================================================================

\begin{table}[h]
\centering
\renewcommand{\arraystretch}{1.25}
\begin{tabularx}{\textwidth}{@{} l X @{}}
\toprule
Symbol & Meaning \\ \midrule
$I,\,i$          & Set and index of known incident locations (calls for service) \\
$J,\,j$          & Set and index of candidate patrol command-centre locations \\
$a_i$            & Weight / priority of incident $i$ (e.g.\ Signal Code value) \\
$P$              & Number of patrol areas (command centres) to be located \\
$S$              & Acceptable service distance (threshold) \\
$d_{ij}$         & Distance from incident $i$ to command centre $j$ \\
$N_i$            & $\{j \in J \mid d_{ij} \le S\}$ --- set of centres that can cover $i$ \\
$x_j$            & 1 if a patrol is located at $j$; 0 otherwise \\
$y_i$            & 1 if incident $i$ is covered by at least one patrol; 0 otherwise \\
$G=(V,E,\ell)$   & OSM directed road graph with edge lengths $\ell(e)$ in metres \\
$v^*_p$          & Nearest OSM node to geographic point $p$ \\
\bottomrule
\end{tabularx}
\caption{Notation common to both approaches.}
\end{table}

% =============================================================================
\section{The Classical PPAC Model (Curtin et al., 2010)}
% =============================================================================

\subsection{Core Objective}

The PPAC model is a direct application of the MCLP.  Its objective is to
select $P$ command-centre locations from the candidate set $J$ so as to
\emph{maximise the total weight of incidents covered} within service
distance~$S$:

\begin{equation}
    \text{Maximise} \quad Z = \sum_{i \in I} a_i \, y_i
    \label{eq:ppac_obj}
\end{equation}

\noindent subject to:
\begin{align}
    \sum_{j \in N_i} x_j &\ge y_i  && \forall\, i \in I
        \label{eq:coverage_constraint}\\
    \sum_{j \in J} x_j   &= P
        \label{eq:cardinality_constraint}\\
    x_j &\in \{0,1\}               && \forall\, j \in J
        \label{eq:binary_x}\\
    y_i &\in \{0,1\}               && \forall\, i \in I
        \label{eq:binary_y}
\end{align}

Constraint~\eqref{eq:coverage_constraint} ensures $y_i=1$ only when at least
one patrol in $N_i$ is selected.  Constraint~\eqref{eq:cardinality_constraint}
fixes the total number of patrols to $P$.

\subsection{Distance Metric}

The paper constructs $N_i$ via an \textbf{all-to-all shortest-path algorithm
on the road network} --- using the North Central Texas Council of Governments
(NCTCOG) street network inside a GIS.  However, the sets $N_i$ are computed
\emph{once and externally}, and the integer programme itself operates on the
binary coverage indicators only.  The optimisation software (CPLEX\,8.1 with
OPL) never sees individual road-link data; it sees only which
$(i,j)$ pairs satisfy $d_{ij} \le S$.

The service distance for Dallas was set to $S = 2$ miles for five of the six
divisions and $S = 1$ mile for the markedly smaller Central Division, chosen
empirically to match typical cross-beat distances.

\subsection{Candidate Facility Locations}

Because no predefined list of admissible command-centre sites existed, the
authors used the \textbf{centroids of the next lower level of the police
hierarchy} (beat centroids within a division) as the candidate set $J$.  This
is acknowledged as a proxy; cadastral data on vacant sites would be preferable.

\subsection{Incident Weights}

Incident weights $a_i$ correspond to the Dallas Police Department's
\emph{Signal Code} priority variable (1--5).  The paper notes that the
relationship between Signal Code and true severity is non-linear, and that
expert input is required to calibrate the weight function properly.

\subsection{Backup Coverage Variant}

The paper introduces a backup (multiple) coverage extension by relaxing
constraint~\eqref{eq:binary_y}:

\begin{equation}
    y_i \in \{0,1,\ldots,P\} \quad \forall\, i \in I
    \label{eq:backup}
\end{equation}

This allows $y_i$ to accumulate each time an additional command centre covers
incident~$i$, rewarding geographic overlap.  A multiobjective formulation is
then solved for a range of values of a coverage lower bound $O$:

\begin{equation}
    \sum_{i=1}^{n} a_i w_i \ge O
    \label{eq:backup_constraint}
\end{equation}

where $w_i$ mirrors the original binary $y_i$, generating a Pareto frontier
between maximal coverage and maximal backup coverage (Fig.\,6 in the paper).

\subsection{Solution Procedure}

The PPAC instances are solved to \textbf{proven optimality} using CPLEX\,8.1
integer programming with branch-and-bound over LP relaxations.  The largest
instance tested involved 87,603 calls for service and was solved in under one
hour on a 2.13\,GHz workstation.

% =============================================================================
\section{Road-Network Distance Extension}
% =============================================================================

\subsection{Motivation}

The PPAC model is explicit that ``any appropriate impedance value (or values)
[can] be used in the computation of the `distances' between incidents and
facilities'' \cite{curtin2010}.  In practice, straight-line distance
systematically \emph{underestimates} actual travel distance in cities with
irregular grids, freeways, water barriers, or hillside terrain --- all of
which are characteristic of Los Angeles.  The road-network extension replaces
the GIS-computed OD matrix with one derived directly from the OSM drivable
graph, making the method fully open-source and reproducible without proprietary
street data.

\subsection{Two-Level Clustering for Candidate Sites}

Rather than using pre-existing police beat centroids as the candidate set $J$,
the extended approach derives candidate locations \emph{from the crime data
itself} using a two-stage weighted K-Means procedure, which closely mirrors the
PPAC objective of concentrating coverage on weighted incidents.

\paragraph{Stage 1 --- Beat generation.}
Minimise the weighted within-cluster sum of squares (WCSS) over $k_1$ clusters:

\begin{equation}
    \underset{\mathcal{S}}{\arg\min}
    \sum_{j=1}^{k_1} \sum_{x_i \in \mathcal{S}_j}
    a_i \,\| x_i - \mu_j \|^2
    \label{eq:wcss}
\end{equation}

where $\mu_j$ is the weighted centroid of beat cluster $\mathcal{S}_j$:

\begin{equation}
    \mu_j = \frac{\displaystyle\sum_{x_i \in \mathcal{S}_j} a_i\, x_i}
                 {\displaystyle\sum_{x_i \in \mathcal{S}_j} a_i}
\end{equation}

\paragraph{Stage 2 --- Sector headquarters.}
The $k_1$ beat centroids are clustered a second time (uniform weights,
$P$ clusters) to yield sector command-centre locations $H = \{h_1,\dots,h_P\}$,
which become the candidate set $J$ for the PPAC evaluation.

\subsection{Road-Network Graph Construction}

The drivable road network for the study area is obtained from OSM via the
\textit{OSMnx} library \cite{boeing2017}:

\begin{equation}
    G = (V, E, \ell), \quad \ell(e) \in \mathbb{R}_{>0} \text{ (metres)}
\end{equation}

The graph covers the convex hull of the city boundary polygon and is persisted
as GraphML for reproducibility.

\subsection{Node Snapping}

Command centres $h_j \in H$ and crime incidents $i \in I$ are geographic
points that may not coincide with an OSM intersection.  Each is snapped to its
nearest graph node:

\begin{equation}
    v^*_p = \underset{v \in V}{\arg\min}\; d_{\mathrm{Euclidean}}(p,\, v)
\end{equation}

This ensures that every distance query is computed between valid road-graph
nodes, avoiding the ``floating point in space'' problem inherent in
straight-line calculations.

\subsection{Road-Network Distance and Coverage}

Shortest-path lengths are computed via Dijkstra's algorithm:

\begin{equation}
    d^{\mathrm{road}}_{ij} =
    \text{Dijkstra}\!\left(G,\; v^*_i,\; v^*_j\right)
    \label{eq:dijkstra}
\end{equation}

The PPAC coverage test then becomes:

\begin{equation}
    N_i = \bigl\{\, j \in J \;\big|\; d^{\mathrm{road}}_{ij} \le S \,\bigr\}
    \label{eq:Ni_road}
\end{equation}

and the objective function and constraints \eqref{eq:ppac_obj}--\eqref{eq:binary_y}
are evaluated unchanged with $N_i$ as defined in \eqref{eq:Ni_road}.

% =============================================================================
\section{Side-by-Side Comparison}
% =============================================================================

\begin{table}[h]
\centering
\renewcommand{\arraystretch}{1.4}
\begin{tabularx}{\textwidth}{@{} l X X @{}}
\toprule
\textbf{Aspect}
    & \textbf{\color{ppacblue}Classical PPAC}
      \textcolor{ppacblue}{(Curtin et al., 2010)}
    & \textbf{\color{roadgreen}Road-Network Extension}
      \textcolor{roadgreen}{(This Work)} \\
\midrule

Study area
    & Dallas, TX (6 police divisions)
    & Los Angeles, CA \\

\addlinespace
Incident data
    & Dallas PD calls for service, 2000\\  (1 day to 1 year, up to 87,603 incidents)
    & LAPD crime incident CSV with \texttt{crime\_weight} field \\

\addlinespace
Candidate locations $J$
    & Centroids of existing beat polygons within each division
    & Derived from data: weighted K-Means beat centroids ($k_1 = 300$),
      then sector clustering ($P = 30$) \\

\addlinespace
Incident weights $a_i$
    & Dallas Signal Code (1--5); non-linear scale acknowledged
    & \texttt{crime\_weight} column from pre-processed dataset \\

\addlinespace
Distance metric $d_{ij}$
    & Shortest path on proprietary NCTCOG road network (inside GIS);
      result exported as OD matrix for the solver
    & Dijkstra shortest path on open OSM drive graph via OSMnx;
      nodes snapped before querying \\

\addlinespace
How $N_i$ is built
    & GIS all-to-all shortest-path query; custom selection interface
    & \texttt{nx.shortest\_path\_length} with \texttt{weight=`length'};
      vectorised over all incident--HQ pairs \\

\addlinespace
Optimisation engine
    & CPLEX 8.1 integer programming (branch-and-bound + simplex LP relaxation);
      \textbf{proven optimal}
    & PPAC coverage metric evaluated post-clustering;
      K-Means heuristic (not globally optimal) \\

\addlinespace
Service distance $S$
    & 2 mi (five divisions); 1 mi (Central Division); empirically chosen
    & 2 mi = 3,218.7 m; configurable parameter \\

\addlinespace
Backup coverage
    & Formally modelled: relaxed $y_i \in \{0,\dots,P\}$;
      Pareto frontier generated via constraint method
    & Not yet implemented; planned as future extension \\

\addlinespace
Geography level solved
    & Sector arrangement \emph{within} each division independently
    & City-wide simultaneously (single clustering pass) \\

\addlinespace
Geometry output
    & Optimal sector boundaries derived from beat assignments
    & Voronoi beat polygons clipped to city boundary; sectors by dissolve \\

\addlinespace
Street data source
    & Proprietary NCTCOG GIS network (Dallas)
    & OpenStreetMap via OSMnx (open, globally available) \\

\addlinespace
Solution quality guarantee
    & \textbf{Optimal} for each division instance
    & Heuristic (K-Means convergence to local minimum) \\

\addlinespace
Computational cost
    & CPLEX solves 87,603-incident instance in $<$ 1 hour
    & K-Means: fast; Dijkstra per incident--HQ pair: $O(|V|\log|V| + |E|)$ each \\

\bottomrule
\end{tabularx}
\caption{Detailed comparison of the classical PPAC model and the road-network extension.}
\label{tab:comparison}
\end{table}

% =============================================================================
\section{Algorithmic Pseudocode}
% =============================================================================

\subsection{Classical PPAC (Curtin et al., 2010)}

\begin{algorithm}[H]
\caption{Classical PPAC --- Maximal Covering for Police Patrol Areas}
\begin{algorithmic}[1]
\State \textbf{Input:} Incidents $I$ (geocoded, weighted $a_i$); candidate sites $J$
       (beat centroids); service distance $S$; patrol count $P$.
\State \textbf{GIS pre-processing:}
    \Statex \quad Compute all-to-all shortest-path OD matrix on the road network.
    \Statex \quad For each $i \in I$: $N_i \leftarrow \{j \in J \mid d_{ij} \le S\}$.
\State \textbf{Export} to integer programming solver:
       $\{|I|,\, |J|,\, a_i,\, N_i\}$.
\State \textbf{Solve PPAC} (CPLEX branch-and-bound):
    \begin{align*}
        &\text{Maximise}\;\; \textstyle\sum_{i \in I} a_i y_i\\
        &\text{s.t.}\;\;
            \textstyle\sum_{j \in N_i} x_j \ge y_i \;\; \forall i;\quad
            \textstyle\sum_j x_j = P;\quad
            x_j, y_i \in \{0,1\}
    \end{align*}
\State \textbf{Assign} beats to sectors based on which optimal facility serves them.
\State \textbf{Evaluate} efficiency: total distance, \% calls covered within $S$.
\State \textbf{Optional --- Backup coverage:}
    \Statex \quad Relax $y_i \in \{0,\dots,P\}$; add constraint $\sum a_i w_i \ge O$.
    \Statex \quad Solve repeatedly over a range of $O$ to trace the Pareto frontier.
\end{algorithmic}
\end{algorithm}

\subsection{Road-Network Extension (This Work)}

\begin{algorithm}[H]
\caption{Road-Network PPAC Extension --- Full Pipeline}
\begin{algorithmic}[1]
\State \textbf{Input:} Crime CSV (LAT, LON, \texttt{crime\_weight}); city boundary GeoJSON;
       $k_1$ (beats); $P$ (sectors); $S$ (miles).
\State \textbf{Pre-processing:}
    \Statex \quad Load and project data to EPSG:3857 (metric Cartesian).
    \Statex \quad Spatial join: retain only incidents inside the city boundary.
\State \textbf{Stage 1 --- Weighted K-Means} ($k_1$ clusters on incidents):
    \Repeat
        \State Assign each incident $i$ to nearest centroid $\mu_j$.
        \State Update $\mu_j \leftarrow \frac{\sum_{x_i \in \mathcal{S}_j} a_i x_i}
                                              {\sum_{x_i \in \mathcal{S}_j} a_i}$.
    \Until{convergence}
    \State Beat centroids $B \leftarrow \{\mu_1, \dots, \mu_{k_1}\}$.
\State \textbf{Stage 2 --- Sector K-Means} ($P$ clusters on $B$):
    \State $H \leftarrow $ K-Means$(B,\, P)$; assign each beat to its nearest sector.
\State \textbf{Stage 3a --- OSM road graph:}
    \Statex \quad Download (or load cached) OSM drive graph $G$ for study area polygon.
\State \textbf{Stage 3b --- Node snapping:}
    \Statex \quad $v^*_j \leftarrow $ nearest node in $G$ for each sector HQ $h_j \in H$.
    \Statex \quad $v^*_i \leftarrow $ nearest node in $G$ for each incident $i \in I$.
\State \textbf{Stage 3c --- PPAC coverage evaluation:}
    \For{each incident $i \in I$}
        \State $d^{\mathrm{road}}_i \leftarrow
               \text{Dijkstra}(G,\; v^*_i,\; v^*_{j^*(i)})$
               \quad where $j^*(i)$ is $i$'s assigned sector.
        \State $y_i \leftarrow \mathbf{1}[d^{\mathrm{road}}_i \le S \cdot 1609.34]$
    \EndFor
    \State $Z \leftarrow \sum_{i \in I} a_i y_i$ \quad (PPAC objective value, Eq.\,\ref{eq:ppac_obj})
\State \textbf{Stage 4 --- Visualisation:}
    \Statex \quad Voronoi tessellation of beat centroids $\to$ clip to city boundary.
    \Statex \quad Dissolve by sector assignment $\to$ sector polygons.
    \Statex \quad Render: sector colours, beat outlines, HQ markers, basemap, coverage stats.
\State \textbf{Output:} PNG patrol map; CSV coverage summary ($Z$, \% covered).
\end{algorithmic}
\end{algorithm}

% =============================================================================
\section{Critical Differences and Limitations}
% =============================================================================

\subsection{Optimality vs.\ Heuristic Quality}

The most fundamental difference is solution quality.  The classical PPAC
produces \textbf{proven optimal} solutions via integer programming: no other
placement of $P$ command centres within the candidate set $J$ can yield a
higher value of $Z$.  The road-network extension uses K-Means, which converges
to a local minimum of the WCSS objective \eqref{eq:wcss} --- not the global
maximum of the PPAC objective \eqref{eq:ppac_obj}.  These are related but
distinct objectives: WCSS minimises average squared distance, while PPAC
maximises weighted coverage within a hard distance threshold.

A natural improvement would be to feed the K-Means sector headquarters into a
proper integer programme (with road-network $N_i$ sets), recovering the
optimality guarantee of the original paper.

\subsection{Scope of the Distance Computation}

In the classical PPAC, road-network distances are used to build $N_i$ but the
integer programme itself is a combinatorial model operating on binary variables.
In the extension, Dijkstra distances are computed for every incident--sector
pair after assignment, meaning the road network is used for \emph{evaluation}
but not for \emph{optimisation}.  The clustering step (which determines command
centre placement) still uses Euclidean distance internally.

\subsection{Geographic Scale}

Curtin et al.\ solve each of Dallas's six divisions \emph{independently},
respecting the administrative hierarchy.  The extension solves the entire LA
city area in a single pass, which is computationally simpler but may produce
sector boundaries that cut across natural administrative divisions.

\subsection{Data Openness}

The classical PPAC relies on proprietary NCTCOG street data and commercial
CPLEX software.  The extension uses exclusively open data (OSM via OSMnx) and
open-source solvers, making it directly reproducible and adaptable to any city
worldwide with OSM coverage.

% =============================================================================
\section{Conclusion}
% =============================================================================

The classical PPAC model of Curtin et al.\ (2010) established that operations
research methods can substantially improve police patrol geography.  The
road-network extension presented here preserves the core PPAC objective but
adapts the method for an open, reproducible workflow.  It replaces proprietary
street data with OSM, derives candidate sites algorithmically from crime data
rather than from a pre-existing police hierarchy, and evaluates coverage using
true drivable distances via Dijkstra shortest paths.

The principal limitation of the extension --- the use of a heuristic (K-Means)
rather than a guaranteed-optimal integer programme --- is also the clearest
direction for future work: coupling the clustering-derived candidate sites with
a proper MCLP solver (e.g., \textit{PuLP}, \textit{OR-Tools}, or an open
CPLEX alternative) would recover the optimality guarantees of the original
paper while retaining the open-data and open-source character of the pipeline.

% =============================================================================
\begin{thebibliography}{9}

\bibitem{curtin2010}
Curtin, K.M., Hayslett-McCall, K., \& Qiu, F. (2010).
\textit{Determining Optimal Police Patrol Areas with Maximal Covering and
Backup Covering Location Models.}
Networks and Spatial Economics, 10, 125--145.

\bibitem{church1974}
Church, R.L., \& ReVelle, C.S. (1974).
\textit{The Maximal Covering Location Problem.}
Papers of the Regional Science Association, 32, 101--118.

\bibitem{hogan1986}
Hogan, K., \& ReVelle, C. (1986).
\textit{Concepts and Applications of Backup Coverage.}
Management Science, 32, 1434--1444.

\bibitem{daskin1981}
Daskin, M.S., \& Stern, E.H. (1981).
\textit{A Hierarchical Objective Set Covering Model for Emergency Medical
Service Vehicle Deployment.}
Transportation Science, 15, 137--152.

\bibitem{boeing2017}
Boeing, G. (2017).
\textit{OSMnx: New Methods for Acquiring, Constructing, Analyzing, and
Visualizing Complex Street Networks.}
Computers, Environment and Urban Systems, 65, 126--139.

\end{thebibliography}

\end{document}